%%%%%%%%%%%%%%%%%%%%%%%%%%%%%%%%%%%%%%%%%%%%%%%%%%%%%%%%%%%%%%%%%%%%%%%%%%%%%%%
% Colo(u)rs of the standard elements
%%%%%%%%%%%%%%%%%%%%%%%%%%%%%%%%%%%%%%%%%%%%%%%%%%%%%%%%%%%%%%%%%%%%%%%%%%%%%%%

%% Individual colours can be modified by the usual \colorlet syntax
%% We also provide a basic interface to bulk-change the 'coupled colours'

\colorlet {ThemeColor}	{white!60!black}	% A general colour which many elements inherit from.
\colorlet {TipColor}	{ThemeColor}		% Tip background
\colorlet {SidebarColor}{ThemeColor}		% Sidebar background
\colorlet {TableColor}  {ThemeColor}		% Table even row background

% Assign a new value to all the co-varying elements
%%If you don't want to co-vary all elements at once, manually use colorlet on the elements to change color
\NewDocumentCommand {\RpgSetThemeColor} { m }
{
	\colorlet {ThemeColor}  {#1}
	\colorlet {TipColor}	{#1}
	\colorlet {SidebarColor}{#1}
	\colorlet {TableColor}  {#1}
}

\colorlet	{PaperColor}		{white}
\definecolor{FiligreeBackColor}	{rgb}	{1,1,1}	 % Background colour of a filigree box
\definecolor{FiligreeColor}		{rgb}	{0,0,0}	 % Colour used to draw the filigree
\definecolor{NarrationColor}	{HTML}	{F7F2E5} % Narration background
\colorlet	{ContourOuterColor} {white!50!black} % Colour of the outer layer of a contoured text
\colorlet	{ContourInnerColor} {black}			 % Colour of the inner core of a contoured text
\colorlet	{SecretColor}		{black!30!white}
